\section*{Appendices: Ready-to-Use Field Instruments}

\subsection*{Appendix A: Channel Experiment Protocol Template}

Before launching any customer acquisition campaign, the founder completes this standardized operational protocol:
\begin{enumerate}[leftmargin=*,itemsep=1pt]
  \item \textbf{Target Buyer and Operational Trigger:} Specific engineering or managerial role; the precise technical or business crisis prompting active exploration.
  \item \textbf{Incumbent Alternatives:} Existing tools, legacy training vendors, or in-house workarounds currently used.
  \item \textbf{Information Gathering Channels:} Specialized technical communities, academic repositories, search terms, or peer networks consulted by the buyer.
  \item \textbf{Primary Acquisition Route:} The single direct channel selected for this experimental test.
  \item \textbf{Supporting Credibility Artifact:} The runnable code repository, technical failure analysis, or evaluation rubric deployed to verify competence.
  \item \textbf{Commercial Offer and Pricing Structure:} Scope of engagement (e.g. 2-weekend practicum, paid architectural diagnostic) and net fee.
  \item \textbf{Resource Expenditure Ceiling:} Absolute financial budget cap (in cash) and maximum allocated founder hours.
  \item \textbf{Target Behavioral Action:} The explicit, low-friction qualification step requested (e.g. submitting an evaluation diagnostic task).
  \item \textbf{Available Delivery Capacity:} Maximum number of students or client engagements that can be served without violating quality thresholds.
  \item \textbf{Observation Window:} Exact start and conclusion dates (typically 14 to 30 days).
  \item \textbf{Disconfirming Empirical Evidence:} Observable data points that will mandate immediate channel termination.
  \item \textbf{Governance Decision Gate:} Continue with current parameters, narrow target audience, pivot positioning, or permanently terminate channel.
\end{enumerate}

\subsection*{Appendix B: Channel Economics and Step-Cost Worksheet}

Let:
\begin{align*}
  P &= \text{Realized net tuition per customer} \\
  V &= \text{Direct variable delivery cost} \\
  A &= \text{Attributable cash acquisition cost} \\
  H &= \text{Founder acquisition time cost}
\end{align*}

The basic unit contribution after customer acquisition is defined as:
\begin{equation}
  C_{\text{unit}} = P - V - A - H
\end{equation}

When customer volume crosses an operational capacity threshold, the organization incurs non-linear \textbf{step costs} (e.g. onboarding an additional qualified code reviewer or leasing dedicated laboratory servers). In these regimes, economic analysis must be conducted at the cohort level:
\begin{equation}
  \begin{aligned}
    \text{Cohort Contribution} &= N \cdot (P - V) - \text{Cash CAC}_{\text{total}} \\
    &\quad - \text{Time CAC}_{\text{total}} - S_{\text{step}}
  \end{aligned}
\end{equation}
where $N$ is the number of enrolled students and $S_{\text{step}}$ represents the incremental fixed staffing step cost.

\subsubsection*{Worked Empirical Comparison}
Consider 2 alternative acquisition motions evaluated by Bayes Institute:
\begin{itemize}[leftmargin=*,itemsep=1pt]
  \item \textbf{Motion 1 (Targeted Direct Outreach):} Acquires 5 learners utilizing ₹10,000 in cash software tools and 20 founder hours valued at an internal planning rate of ₹800 per hour (₹16,000 founder time cost).
  \begin{equation*}
    \text{Economic CAC}_1 = \frac{\text{₹}10,000 + \text{₹}16,000}{5} = \text{₹}5,200
  \end{equation*}
  \item \textbf{Motion 2 (Paid Practitioner Search Ads):} Acquires 5 learners utilizing ₹40,000 in cash advertising spend and 5 founder hours valued at ₹800 per hour (₹4,000 founder time cost).
  \begin{equation*}
    \text{Economic CAC}_2 = \frac{\text{₹}40,000 + \text{₹}4,000}{5} = \text{₹}8,800
  \end{equation*}
\end{itemize}

While Motion 1 is significantly cheaper on an economic CAC basis (₹5,200 vs ₹8,800), Motion 1 consumes 20 founder hours for every 5 learners, capping monthly acquisition at 10 students. If Motion 2 can reliably deliver 25 qualified learners per month at unchanged marginal economics and instructional capacity exists, Motion 2 becomes the superior scaling engine despite its higher initial CAC. Both assumptions must be tested empirically.

\subsection*{Appendix C: Buyer-Journey Investigation Guide}

To uncover authentic buying behavior without superficial market research bias, the founder interviews a recently enrolled customer or qualified prospect, reconstructing their most recent comparable purchase:
\begin{itemize}[leftmargin=*,itemsep=1pt]
  \item \textbf{The Catalyst Event:} What specific incident, system outage, or project requirement prompted you to look for a solution?
  \item \textbf{Initial Search Behavior:} Exactly what search queries did you enter into Google or GitHub? Who was the 1st colleague you contacted?
  \item \textbf{Evaluation Consideration Set:} Which competing alternatives, courses, or consultancies did you evaluate? Why were others eliminated?
  \item \textbf{Credibility Thresholds:} What specific piece of documentation, code, or interaction convinced you that this offering was legitimate?
  \item \textbf{Internal Approval Rail:} Did you pay out of personal discretionary funds, or secure corporate departmental reimbursement? What procurement steps were required?
  \item \textbf{Decision Latency:} How many days elapsed between your 1st website visit and your final payment submission? What caused pauses or delays?
  \item \textbf{Touchpoint Audit:} Which emails, articles, or videos were genuinely helpful, and which were completely ignored?
\end{itemize}

Founders must never ask abstract questions such as ``Which social networks do you enjoy?'' Personal media preferences do not reflect professional purchasing behavior. Ground every interview question in concrete, past behavioral facts.

\subsection*{Appendix D: Partnership Activation Agreement Checklist}

Before signing any formal distribution or channel partnership agreement, verify that the following 12 operational clauses are documented:
\begin{enumerate}[leftmargin=*,itemsep=1pt]
  \item \textbf{Partner Economic Rationale:} The documented commercial or strategic value the partner derives from the relationship.
  \item \textbf{Exact Target Audience Definition:} Explicit professional roles, membership tiers, or corporate departments to be contacted.
  \item \textbf{Permitted Communication Channels:} Pre-agreed member newsletters, portal catalog listings, or dedicated co-branded events.
  \item \textbf{Operational Ownership:} Specific named executives on both sides assigned direct operational execution responsibility.
  \item \textbf{Commercial Economics and Revenue Share:} Transparent commission splits (e.g. 15\% net), invoicing rails, and payment schedules.
  \item \textbf{Legal Referral Disclosures:} Mandatory ethical disclosures informing prospective buyers of the partner's financial interest.
  \item \textbf{Trademark and Intellectual Property Rights:} Strict boundaries governing brand representation, logo usage, and academic claims.
  \item \textbf{Data Protection and Privacy Governance:} Full compliance with applicable statutory data-sharing and consent regulations.
  \item \textbf{Candidate Qualification Standards:} Explicit admissions prerequisite screening protocols to prevent unqualified enrollments.
  \item \textbf{Service and Support Allocations:} Clear division of customer onboarding, technical infrastructure, and dispute resolution duties.
  \item \textbf{Termination and Exit Protocols:} Unilateral termination rights if agreed quality, ethical, or commercial benchmarks are breached.
  \item \textbf{First Measurable Activation Milestone:} An explicit 30-day operational deliverable (e.g. hosting 1 joint technical seminar).
\end{enumerate}
